% Generated draft: requires mathematical and source review.
\documentclass[11pt]{article}
\usepackage[margin=25mm]{geometry}
\usepackage{fontspec}
\setmainfont{Latin Modern Roman}
\usepackage{amsmath,amssymb,mathtools,unicode-math}
\setmathfont{Latin Modern Math}
\usepackage{enumitem,needspace,xurl,hyperref}
\setlength{\parindent}{0pt}
\setlength{\parskip}{5pt}
\setlength{\emergencystretch}{3em}
\newcommand{\N}{\mathbb{N}}
\newcommand{\Z}{\mathbb{Z}}
\newcommand{\Q}{\mathbb{Q}}
\newcommand{\R}{\mathbb{R}}
\newcommand{\C}{\mathbb{C}}
\newcommand{\F}{\mathbb{F}}
\newcommand{\A}{\mathbb{A}}
\newcommand{\PP}{\mathbb{P}}
\newcommand{\E}{\mathbb{E}}
\newcommand{\eps}{\varepsilon}
\newcommand{\dd}{\,\mathrm{d}}
\begin{document}
\section*{UnsolvedMath ID: 30002785}
\textbf{Source problem code:} catalogue entry
\subsection{Definitions and mathematical statement}
The supplied record asks two questions about operator ideals denoted by \(L_q^{\mathrm{hilb}}\), \(L_p^{\mathrm{app}}\), and \(L_p^{\mathrm{ent}}\). It does not define these ideals, specify their underlying number sequences or membership conditions, give the admissible exponent range, or identify the intended class of spaces. These omissions prevent a fully determined mathematical statement.

For clarity, the standard Banach-space meaning of approximability is as follows. For Banach spaces \(E,F\) over a fixed scalar field, let \(\mathcal L(E,F)\) be the bounded linear operators and let \(\mathcal F(E,F)\) be those with finite-dimensional range. Define
\[
\mathcal A(E,F)=\overline{\mathcal F(E,F)}^{\|\cdot\|}.
\]
Thus \(T\) is approximable precisely when, for every \(\varepsilon>0\), there exists \(S\in\mathcal F(E,F)\) with \(\|T-S\|<\varepsilon\). This definition states a conventional interpretation; the supplied record does not confirm the ambient Banach-space setting, scalar field, or intended meaning of approximability.

To retain the source quantifiers without supplying missing hypotheses, let \(D\) stand for the unspecified admissible exponent set and let \(\mathscr C\) stand for the unspecified class of ordered pairs of spaces. These are placeholders requiring source clarification. With the three ideal symbols retaining their source meanings, the first question is whether
\[
\forall p,q\in D,\qquad
\frac1q\geq\frac1p+1
\quad\Longrightarrow\quad
\forall(E,F)\in\mathscr C,\quad
L_q^{\mathrm{hilb}}(E,F)\subseteq L_p^{\mathrm{app}}(E,F).
\]
The equality case \(1/q=1/p+1\) is explicitly included. Any convention for infinite exponents must be specified before interpreting this formula at such endpoints.

The second question is to determine all admissible \(p\) such that every operator in \(L_p^{\mathrm{ent}}\) is approximable. Under the Banach-space interpretation above, its schematic target is the set
\[
\left\{p\in D:\ 
\forall(E,F)\in\mathscr C\ \forall T\in L_p^{\mathrm{ent}}(E,F),\quad
T\in\mathcal A(E,F)\right\}.
\]
Neither formula resolves the missing definitions or hypotheses. In particular, the superscripts alone do not specify the operator ideals or their sequence summability requirements.
\subsection{Short English statement}
The source asks whether \(L_q^{\mathrm{hilb}}\subseteq L_p^{\mathrm{app}}\) whenever \(1/q\geq1/p+1\), including equality, and which exponents \(p\) make every operator in \(L_p^{\mathrm{ent}}\) approximable. The ideal definitions, exponent ranges, and ambient spaces require clarification. The displayed statements preserve these questions schematically and require human mathematical review.
\subsection{Sources}
\begin{itemize}
\item Local UnsolvedMath registry: supplied record with permanent ID \(30002785\), title \emph{Inclusions and Approximability in Operator Ideals}, source problem code \texttt{catalogue entry}, in the Oberwolfach Reports Open Problems collection. The record gives proposed year 2015 and identifies no individual proposer.
\item Supplied source citation: \emph{Mini-Workshop: Modern Applications of \(s\)-numbers and Operator Ideals} (2015), Oberwolfach Reports, pp.\ 369--397, DOI \texttt{10.4171/owr/2015/6}. Supplied source URL and literature reference labeled ``original OWR report'': \url{https://doi.org/10.4171/owr/2015/6}. The supplied background describes Oberwolfach Reports as documenting mathematical workshops and their open-problem discussions. The report contents were not independently inspected for this draft.
\item Supplied literature reference labeled ``recent inequalities between s-numbers'': \url{https://doi.org/10.1007/s43036-024-00386-x}. No further bibliographic details are supplied.
\item Literature assessment checked at \texttt{2026-08-21}: supplied dated curation, not independently verified. It labels the record open and describes it as a collection of operator-ideal questions, including unresolved inclusion and approximability problems; it states that standard inequalities between \(s\)-numbers do not settle the selected questions. The supplied background calls this dated literature triage derived from the OWR curation and requests independent verification of time-sensitive status and cited claims.
\item The registry labels the statement \texttt{corrected\_verified} and says its mathematical content was preserved while notation, typography, or omitted definitions were repaired for publication. This is supplied registry metadata, not independent verification of the formulation or mathematical status in this draft.
\end{itemize}
\end{document}
